\section{Lower bound on the maximum incorrect logit (anti-concentration)}
\label{sec:incorrect_max_lower}

The failure branch of \Cref{thm:R1} needs the opposite of
\Cref{lem:incorrect_max}: a high-probability \emph{lower} bound on the
maximum incorrect logit, showing that when the descent rate is
sub-critical \emph{some} incorrect token overtakes the correct one. This
is a direct anti-concentration statement on the $|\Vocab|-1$ incorrect
logits, obtained from the two classical \emph{lower}-bound tools for the
supremum of a Gaussian process --- Sudakov minoration for the expected
maximum and the Borell--TIS Gaussian concentration inequality to pin the
maximum to its mean. It is \emph{not} a reversed version of
\Cref{lem:incorrect_max} (a union bound can only \emph{upper}-bound a max)
and \emph{not} a branching-process coupling.

\begin{lemma}[Maximum incorrect logit, lower bound]
\label{lem:incorrect_max_lower}
Suppose \Cref{ass:descent,ass:bounded_architecture} hold with
$|\Vocab|\ge3$, and fix a horizon $T\ge1$; in particular invoke the
isotropic-Gaussian noise form \Cref{eq:noise_gaussian} of
\Cref{ass:bounded_architecture}(5) (used by this lemma only,
\Cref{rem:gaussian_noise_scope}). Let
$\sigma_T\asymp R_U M\,e^S/\sqrt{T d}$ be the per-direction fluctuation
scale of the incorrect logits (the same scale as in
\Cref{lem:incorrect_max,lem:azuma}). Then there are absolute constants
$\kappa\in(0,1)$ and $c_3>0$ such that, on an event of probability at least
$1-(|\Vocab|-1)^{-c_3}$,
\begin{align}\label{eq:incorrect_max_lower_bound}
   \max_{a\neq\corr}\, \inner{W_U^a}{x_T}
   \;\ge\;
   \sigma_T\,\sqrt{2\,\kappa\,(1-\incoh_0)\,\log(|\Vocab|-1)}
   \;-\;
   \incoh_0\, R_U\, \abs D,
\end{align}
where $\abs D=\abs{\inner{\rowhat}{\E\tilde x_T}}\le\snet M$ is the shared
correct-direction drift magnitude of \Cref{lem:incorrect_max} (so the
deterministic correction is $O(\critrate)$ in the sub-critical regime, not
$\Theta(1)$).
\end{lemma}

\begin{proof}
Under the Gaussian noise form \Cref{eq:noise_gaussian} the centred
incorrect projections are an exactly Gaussian vector, and the bound follows
from two standard lower-bound tools for the supremum of a Gaussian process:
Sudakov minoration (\cite{vershynin2018}, Theorem~7.4.1) lower-bounds the
\emph{expected} maximum through the pairwise separation supplied by the
incoherence, and the Borell--TIS Gaussian concentration inequality
(\cite{vershynin2018}, Theorem~5.2.3) pins the maximum to that expectation.
This is a genuine lower bound: a union bound can only \emph{upper}-bound a
max (\Cref{lem:incorrect_max}); the lower bound needs Gaussian comparison
plus the pairwise decorrelation of the incoherence, and there is no
Galton--Watson coupling.

\smallskip\noindent\textbf{Step 0 (the centred projections form a Gaussian
vector).} For $a\neq\corr$ write the centred incorrect projection as the
noise aggregate
\begin{align*}
   Z_a \;\coloneqq\; \inner{W_U^a}{\tilde x_T-\E\tilde x_T}
       \;=\; \sum_{k=0}^{T} w_{T,k}\,\inner{W_U^a}{\xi_k},
\end{align*}
using $\tilde x_T=\sum_k w_{T,k}V_k$ (\Cref{eq:running_average_teaser}) and
the drift/noise split $V_k=\E[V_k\mid\Fcal_{k-1}]+\xi_k$ of
\Cref{ass:bounded_architecture}(5). Let $\Gcal$ be the $\sigma$-field
generated by the attention weights $\{w_{T,k}\}_{k}$ and the conditional
means $\{\E[V_k\mid\Fcal_{k-1}]\}_k$ (the predictable data of the
trajectory), so that under $\Gcal$ the weights $w_{T,k}$ and the scalars
$\tau_k^2$ of \Cref{eq:noise_gaussian} are fixed and the only remaining
randomness is the noise $(\xi_k)_k$. By \Cref{eq:noise_gaussian} each
$\xi_k\mid\Fcal_{k-1}$ is Gaussian, and the $(\xi_k)_k$ are a
martingale-difference sequence; conditioning on $\Gcal$ therefore makes
$(\xi_k)_k$ a jointly Gaussian \emph{independent} family with
$\xi_k\sim\mathcal N(0,\tau_k^2 I_d)$. Hence, conditional on $\Gcal$, the
family $(Z_a)_{a\neq\corr}$ --- a fixed linear image of $(\xi_k)_k$ --- is
a centred Gaussian vector. Cross terms across $k$ vanish by independence,
and the per-step covariance is the \emph{scalar} $\tau_k^2 I_d$ (this is the
one place the scalar covariance of \Cref{eq:noise_gaussian}, rather than
the merely comparable-to-isotropic \Cref{eq:noise_isotropy}, is used), so
\begin{align}\label{eq:antic_gram_cov}
   \operatorname{Cov}(Z_a,Z_{a'}\mid\Gcal)
   \;=\; \sum_{k} w_{T,k}^2\,\tau_k^2\,\inner{W_U^a}{W_U^{a'}}
   \;=\; \beta^2\,\inner{W_U^a}{W_U^{a'}},
   \qquad
   \beta^2\coloneqq\sum_k w_{T,k}^2\,\tau_k^2 .
\end{align}

\smallskip\noindent\textbf{Step 1 (variance scale, exact correlation, and
pairwise separation).} The common scalar $\beta^2$ sets the fluctuation
scale. Since $\sigma_{\min}^2/d\le\tau_k^2\le\sigma_{\max}^2/d$,
\begin{align}\label{eq:antic_beta}
   \tfrac{\sigma_{\min}^2}{d}\,S_T \;\le\; \beta^2 \;\le\; \tfrac{\sigma_{\max}^2}{d}\,S_T,
   \qquad
   S_T\coloneqq\sum_k w_{T,k}^2 \;\in\; \bigl[\,\tfrac1T,\ \tfrac{e^{2S}}{T}\,\bigr],
\end{align}
where the two-sided range of $S_T$ is the trivial floor $S_T\ge(\sum_k w_{T,k})^2/T=1/T$
(Cauchy--Schwarz with $\sum_k w_{T,k}=1$) together with the bounded-score
upper bound $S_T\le e^{2S}/T$ (\Cref{lem:quad_variation},
\Cref{eq:ST_bound}). Hence $\beta=\Theta(\sigma_T)$ with the fixed scale
$\sigma_T\asymp R_U M e^S/\sqrt{Td}$ of \Cref{lem:incorrect_max}: both
$\beta\le\sigma_{\max}e^S/\sqrt{Td}$ and
$\beta\ge\sigma_{\min}/\sqrt{Td}$, and the gap between the two is the
$e^S=\Theta(1)$ prefactor of the bounded-score regime
\Cref{ass:bounded_architecture}(4) (with $\sigma_{\min},\sigma_{\max}=\Theta(M)$,
\Cref{rem:incorrect_max_origin}), absorbed into the absolute constant
$\kappa$ below exactly as it is in the upper bound \Cref{lem:incorrect_max}.
Setting $a'=a$ in
\Cref{eq:antic_gram_cov} gives the diagonal
$\operatorname{Var}(Z_a\mid\Gcal)=\beta^2\norm{W_U^a}_2^2\in[\beta^2\rho_0^2,\beta^2 R_U^2]$.
Because the per-step covariance is a scalar multiple of the identity, the
correlation reduces \emph{exactly} to the normalised Gram entry of
\Cref{def:incoherence}: for $a\neq a'$,
\begin{align}\label{eq:antic_corr}
   \abs{\operatorname{Corr}(Z_a,Z_{a'}\mid\Gcal)}
   \;=\; \frac{\abs{\inner{W_U^a}{W_U^{a'}}}}{\norm{W_U^a}_2\,\norm{W_U^{a'}}_2}
   \;\le\; \incoh(W_U) \;\le\; \incoh_0 ,
\end{align}
with no Berry--Esseen correction and \emph{no} restriction on the size of
$\incoh_0$ --- this exactness is precisely what the Gaussian form
\Cref{eq:noise_gaussian} buys (cf.\ \Cref{rem:anti_concentration_matching}).
The \emph{canonical (}$L^2$\emph{) metric} of the Gaussian vector then
separates distinct directions. Directly from \Cref{eq:antic_gram_cov},
\begin{align}\label{eq:antic_delta}
   \norm{Z_a-Z_{a'}}_2^2
   \;=\; \beta^2\,\norm{W_U^a-W_U^{a'}}_2^2
   \;=\; \beta^2\bigl(\norm{W_U^a}_2^2+\norm{W_U^{a'}}_2^2-2\inner{W_U^a}{W_U^{a'}}\bigr),
\end{align}
an exact identity (the canonical metric is the $\beta$-scaled Euclidean
metric of the rows, again because the per-step covariance is scalar). Using
$\norm{W_U^a}_2,\norm{W_U^{a'}}_2\ge\rho_0$ and the Gram bound
$\inner{W_U^a}{W_U^{a'}}\le\incoh(W_U)\norm{W_U^a}_2\norm{W_U^{a'}}_2\le\incoh_0 R_U^2$
(\Cref{def:incoherence,ass:bounded_architecture}(1)) gives
$\norm{Z_a-Z_{a'}}_2^2\ge2\beta^2(\rho_0^2-\incoh_0 R_U^2)$. Taking the
standard normalised convention $\rho_0=R_U$ of
\Cref{rem:architecture_realism} (the general case only rescales the
absolute constant), $\norm{Z_a-Z_{a'}}_2^2\ge2\beta^2\rho_0^2(1-\incoh_0)$,
i.e.
\begin{align}\label{eq:antic_delta_def}
   \norm{Z_a-Z_{a'}}_2
   \;\ge\; \sqrt2\,\beta\rho_0\sqrt{1-\incoh_0}
   \;\ge\; \delta_{\mathrm{sep}}
   \;\coloneqq\; c_0\,\sigma_T\sqrt{2(1-\incoh_0)}
   \qquad\text{for all }a\neq a',
\end{align}
where $c_0\coloneqq\rho_0\inf_{\Gcal}(\beta/\sigma_T)>0$ is an absolute
constant (independent of $\Gcal$) by the two-sided scale
$\beta\ge\sigma_{\min}/\sqrt{Td}=\Omega(\sigma_T)$ of \Cref{eq:antic_beta};
it folds in $\rho_0$ and the $e^{-S}=\Theta(1)$ prefactor.

\smallskip\noindent\textbf{Step 2 (Sudakov minoration: the expected
maximum).} Work conditional on $\Gcal$, where $(Z_a)_{a\neq\corr}$ is a
centred Gaussian process on the index set $S\coloneqq\Vocab\setminus\{\corr\}$,
$|S|=|\Vocab|-1=:m$, with canonical metric $d_Z(a,a')=\norm{Z_a-Z_{a'}}_2$.
By \Cref{eq:antic_delta_def} every pair is at canonical distance
$\ge\delta_{\mathrm{sep}}$, so the open balls of radius
$\varepsilon\coloneqq\delta_{\mathrm{sep}}/2$ around the $m$ points are
disjoint and any $\varepsilon$-net must contain at least one point per
ball: the covering number satisfies $N(S,d_Z,\varepsilon)\ge m$. Sudakov
minoration (\cite{vershynin2018}, Theorem~7.4.1; digest
\texttt{sudakov-minoration.md}) gives, for an absolute constant $c>0$,
\begin{align}\label{eq:antic_sudakov}
   \E\!\bigl[\max_{a\in S}Z_a \;\big|\; \Gcal\bigr]
   \;\ge\; c\,\varepsilon\,\sqrt{\log N(S,d_Z,\varepsilon)}
   \;\ge\; c\,\tfrac{\delta_{\mathrm{sep}}}{2}\,\sqrt{\log m}
   \;\ge\; c_1'\,\sigma_T\,\sqrt{(1-\incoh_0)\,\log(|\Vocab|-1)},
\end{align}
with $c_1'\coloneqq c\,c_0/\sqrt2>0$ absolute, using \Cref{eq:antic_delta_def}
and $m=|\Vocab|-1$.

\smallskip\noindent\textbf{Step 3 (Borell--TIS concentration: pinning the
maximum to its mean).} Still conditional on $\Gcal$, write
$f(\,\cdot\,)\coloneqq\max_{a\in S}(\,\cdot\,)_a$ and represent the Gaussian
vector as $Z=\Sigma^{1/2}_S\,g$ with $g\sim\mathcal N(0,I)$ and
$\Sigma_S=\operatorname{Cov}((Z_a)_{a\in S}\mid\Gcal)$. The map
$g\mapsto f(\Sigma_S^{1/2}g)$ is Lipschitz with constant
$\max_{a\in S}\norm{\Sigma_S^{1/2}e_a}_2=\max_{a\in S}\operatorname{Var}(Z_a\mid\Gcal)^{1/2}\le\beta R_U\le C'\sigma_T$
(by the diagonal $\operatorname{Var}(Z_a\mid\Gcal)=\beta^2\norm{W_U^a}_2^2\le\beta^2 R_U^2$
of Step~1, with $C'\coloneqq R_U\sup_{\Gcal}(\beta/\sigma_T)<\infty$ a
$\Gcal$-independent absolute constant by the upper scale
$\beta\le\sigma_{\max}e^S/\sqrt{Td}=O(\sigma_T)$ of \Cref{eq:antic_beta}),
since
$\abs{\max_a u_a-\max_a v_a}\le\max_a\abs{u_a-v_a}\le\norm{u-v}_2$. The
Borell--TIS / Gaussian concentration inequality (\cite{vershynin2018},
Theorem~5.2.3; digest \texttt{borell-tis.md}) then yields, for every
$u\ge0$,
\begin{align}\label{eq:antic_borelltis}
   \Pr\!\Bigl[\max_{a\in S}Z_a \;<\; \E[\max_{a\in S}Z_a\mid\Gcal]-u \;\Big|\; \Gcal\Bigr]
   \;\le\; \exp\!\Bigl(-\tfrac{c_2\,u^2}{(C'\sigma_T)^2}\Bigr).
\end{align}
Choose $u\coloneqq c_1'\sigma_T\sqrt{(1-\incoh_0)\log(|\Vocab|-1)}/2$, i.e.\
half the Sudakov lower bound \Cref{eq:antic_sudakov}. Then by
\Cref{eq:antic_sudakov},
$\E[\max_a Z_a\mid\Gcal]-u\ge\tfrac12 c_1'\sigma_T\sqrt{(1-\incoh_0)\log(|\Vocab|-1)}$,
and \Cref{eq:antic_borelltis} bounds the failure probability by
$\exp(-c_2(c_1'/(2C'))^2(1-\incoh_0)\log(|\Vocab|-1))=(|\Vocab|-1)^{-c_3}$
with the absolute constant
$c_3\coloneqq c_2(c_1'/(2C'))^2(1-\incoh_0)>0$. Set
$\kappa\coloneqq(c_1'/2)^2/2\in(0,1)$, so that
$\tfrac12 c_1'\sqrt{1-\incoh_0}=\sqrt{2\kappa(1-\incoh_0)}$. Since the bound
\Cref{eq:antic_borelltis} holds conditional on $\Gcal$ with a probability
not depending on $\Gcal$, it holds unconditionally after taking
expectation over $\Gcal$. Hence, on an event of probability at least
$1-(|\Vocab|-1)^{-c_3}$,
\begin{align}\label{eq:antic_maxZ}
   \max_{a\neq\corr}Z_a
   \;=\; \max_{a\neq\corr}\inner{W_U^a}{\tilde x_T-\E\tilde x_T}
   \;\ge\; \sigma_T\,\sqrt{2\,\kappa\,(1-\incoh_0)\,\log(|\Vocab|-1)}.
\end{align}
The constant $\kappa\in(0,1)$ rescales the $\sqrt{1-\incoh_0}$ threshold in
\Cref{eq:incorrect_max_lower_bound} by a $\Theta(1)$ amount, absorbed into
the constant $c_1$ of \Cref{thm:R1}; it carries no dependence on the
numerical values of the Sudakov / Borell--TIS constants $c,c_2$.

\smallskip\noindent\textbf{Step 4 (the shared deterministic drift).} On the
event of \Cref{eq:antic_maxZ}, let $\hat a\neq\corr$ attain the maximum;
adding back the mean,
\begin{align*}
   \inner{W_U^{\hat a}}{\tilde x_T}
   \;=\; Z_{\hat a}+\inner{W_U^{\hat a}}{\E\tilde x_T}
   \;\ge\; \sigma_T\sqrt{2\kappa(1-\incoh_0)\log(|\Vocab|-1)} - \incoh_0 R_U\abs D,
\end{align*}
where $\inner{W_U^{\hat a}}{\E\tilde x_T}=D\inner{W_U^{\hat a}}{\rowhat}\ge-\incoh_0 R_U\abs D$
carries the \emph{shared} correct-direction drift magnitude
$\abs D=\abs{\inner{\rowhat}{\E\tilde x_T}}\le\snet M$ (\Cref{lem:incorrect_max};
in the sub-critical regime $\snet\le\critrate/\sqrt2$ this is
$O(\critrate)$, hence lower-order). Therefore
$\max_{a\neq\corr}\inner{W_U^a}{\tilde x_T}\ge\sigma_T\sqrt{2\kappa(1-\incoh_0)\log(|\Vocab|-1)}-\incoh_0 R_U\abs D$,
which is \Cref{eq:incorrect_max_lower_bound} (with the deterministic term
sharpened from $\incoh_0 R_U M$ to $\incoh_0 R_U\abs D$); the transfer to
$x_T$ via the exact rescaling \Cref{eq:onsphere_direction} adds
$O(1/d^{1.5})$ (\Cref{lem:retraction_stability}), absorbed into $\sigma_T$.

\smallskip\noindent\textbf{Probability budget.} The
$1-(|\Vocab|-1)^{-c_3}$ probability is the \emph{single} Borell--TIS
deviation event \Cref{eq:antic_borelltis} for the functional
$\max_{a\neq\corr}Z_a$; \emph{no} union over the $|\Vocab|-1$ directions is
taken (a union would yield the wrong, upper, direction --- precisely
\Cref{lem:incorrect_max}). The $|\Vocab|-1$ directions enter only through
the metric-entropy count $N(S,d_Z,\varepsilon)\ge|\Vocab|-1$ inside the
Sudakov bound \Cref{eq:antic_sudakov}. This completes the proof.
\end{proof}

\begin{remark}[Why the Gaussian route, and matching the upper bound]
\label{rem:anti_concentration_matching}
\emph{(i) Transfer route.} Sudakov minoration and Borell--TIS are
lower-bound tools for \emph{Gaussian} processes; we therefore invoke the
isotropic-Gaussian form \Cref{eq:noise_gaussian} of
\Cref{ass:bounded_architecture}(5), under which the centred incorrect
projections $(Z_a)_{a\neq\corr}$ are \emph{exactly} jointly Gaussian
(Step~0) and the two tools apply directly, \emph{for every}
$\incoh_0\in[0,1)$, with no regime restriction. This is the honest cost of
an unconditional two-sided transition: a second-moment / Paley--Zygmund
argument would instead need the exceedance covariance bound
$\operatorname{Cov}(\1\{Z_a>t\},\1\{Z_{a'}>t\})\le\incoh_0\Pr\Pr$, which is
\emph{false} for constant $\incoh_0$ at the relevant threshold
$t\asymp\sigma_T\sqrt{(1-\incoh_0)\log|\Vocab|}$ --- there two
$\incoh_0$-correlated Gaussians have
$\Pr[\text{both}>t]/(\Pr\Pr)=|\Vocab|^{\Theta(\incoh_0)}$, so the
pair-covariance sum dominates $(\E N)^2$ and the variance is not
$o((\E N)^2)$ unless $\incoh_0=O(1/\log|\Vocab|)$. The Sudakov/Borell--TIS
route avoids the second moment entirely and so closes the failure branch
unconditionally in $\incoh_0$. \emph{The Gaussian law is used here only;}
the success branch and \Cref{thm:R2} use solely the second-moment bound
\Cref{eq:noise_isotropy} (\Cref{rem:gaussian_noise_scope}).

\smallskip
\emph{(ii) Two-sided transition.} The lower and upper thresholds (the
latter from \Cref{lem:incorrect_max}),
\begin{align*}
   \sigma_T\sqrt{2\kappa(1-\incoh_0)\log(|\Vocab|-1)}
   \qquad\text{vs.}\qquad
   \sigma_T\sqrt{2\log(|\Vocab|-1)},
\end{align*}
both with $\sigma_T\asymp R_U M e^S/\sqrt{T_{\max}d}$, match up to the
explicit constant $\kappa(1-\incoh_0)\le1$. This matching is what makes
\Cref{thm:R1} a genuine \emph{two-sided} phase transition at $\critrate$:
above $\sqrt2\,\critrate$ the signal floor clears the incorrect-max upper
bound (success), and below $\critrate/\sqrt2$ the incorrect-max lower bound
clears the signal ceiling (failure). The lower isotropy bound
$\sigma_{\min}>0$ in \Cref{eq:noise_gaussian} supplies the non-degeneracy
(the canonical-metric separation \Cref{eq:antic_delta_def}); if it failed the
incorrect directions could collapse and no anti-concentration would hold.
\end{remark}
